% =============================================================================
% Résumé exécutif — Projet Data Science sous R (ESME, 2026-2027)
%
% Document de 2 pages au maximum, destiné à un public NON technique
% (direction marketing) : pas de code, pas de jargon statistique, des chiffres
% utiles à la décision.
%
% Mode d'emploi :
%   1. Remplacez chaque [À compléter] et supprimez les commentaires de consigne.
%   2. Exportez depuis R le graphique qui porte votre message principal :
%        ggsave("figures/figure-cle.pdf", width = 7, height = 3.5)
%      Tant que ce fichier n'existe pas, un cadre de remplacement s'affiche.
%   3. Compilez avec pdfLaTeX (RStudio : bouton « Compile PDF » ; ou Overleaf ;
%      ou en ligne de commande : pdflatex resume-executif.tex).
%   4. Vérifiez que le document tient en 2 pages : sinon, un avertissement
%      s'affiche en bas de la dernière page.
% =============================================================================

\documentclass[11pt, a4paper]{article}

\usepackage[utf8]{inputenc}
\usepackage[T1]{fontenc}
\usepackage[french]{babel}
\usepackage{lmodern}
\usepackage{microtype}
\usepackage[margin=2cm, top=1.6cm, bottom=1.8cm]{geometry}
\usepackage[table]{xcolor}
\usepackage{graphicx}
\usepackage{booktabs}
\usepackage{tabularx}
\usepackage{enumitem}
\usepackage[most]{tcolorbox}
\usepackage{titlesec}
\usepackage[hidelinks]{hyperref}

% ----------------------------------------------------------------------------
% Informations à compléter
% ----------------------------------------------------------------------------
\newcommand{\numeroprojet}{X}
\newcommand{\titreprojet}{[À compléter : titre qui énonce la question ou le message principal]}
\newcommand{\membres}{Prénom NOM, Prénom NOM, Prénom NOM, Prénom NOM}
\newcommand{\groupe}{Groupe X}

% ----------------------------------------------------------------------------
% Mise en forme (inutile de modifier)
% ----------------------------------------------------------------------------
\definecolor{principal}{HTML}{1565C0}
\definecolor{fondencadre}{HTML}{E8F1FB}
\setlength{\parindent}{0pt}
\setlength{\parskip}{0.45em}
\setlist{nosep, leftmargin=1.4em}
\titleformat{\section}{\large\bfseries\color{principal}}{}{0em}{}
\titlespacing*{\section}{0pt}{0.9em}{0.3em}
\pagestyle{plain}

\newtcolorbox{enbref}{enhanced, colback=fondencadre, colframe=principal,
  boxrule=0pt, leftrule=4pt, arc=2pt, left=8pt, right=8pt, top=4pt, bottom=5pt,
  title={En bref}, fonttitle=\bfseries, coltitle=principal,
  colbacktitle=fondencadre, toptitle=3pt, bottomtitle=0pt}

% Signale un dépassement de la limite de 2 pages en bas de la dernière page
\AtEndDocument{%
  \ifnum\value{page}>2
    \par\vfill
    \begin{center}\color{red}\bfseries
      Attention : le résumé exécutif dépasse 2 pages (\thepage{} pages).
    \end{center}
    \typeout{ATTENTION : le resume executif depasse 2 pages.}
  \fi}

\begin{document}

% ----------------------------------------------------------------------------
% En-tête
% ----------------------------------------------------------------------------
{\small\color{gray} Data Science sous R — ESME, Ingé3 2026-2027 \hfill
 Projet n°\,\numeroprojet{} — \groupe}

\vspace{0.4em}
{\LARGE\bfseries\color{principal} \titreprojet\par}
\vspace{0.3em}
{\small \membres{} \hfill \today}

\vspace{0.2em}
\rule{\linewidth}{0.4pt}

% ----------------------------------------------------------------------------
\begin{enbref}
% Consigne : les trois messages à retenir si le lecteur ne lit que cet encadré.
% Une phrase chacun, avec un chiffre.
\begin{itemize}
  \item \textbf{[Constat principal]} : [À compléter, avec un chiffre clé].
  \item \textbf{[Deuxième résultat]} : [À compléter].
  \item \textbf{[Recommandation principale]} : [À compléter, avec l'effet attendu].
\end{itemize}
\end{enbref}

% ----------------------------------------------------------------------------
\section{Le contexte et la question}
% Consigne : 3 à 4 phrases. Quel problème marketing ? Quelle décision doit être
% prise ? Quelle question précise avons-nous étudiée ?
[À compléter : situation de l'entreprise, enjeu, question posée et décision
que l'analyse doit éclairer.]

% ----------------------------------------------------------------------------
\section{Ce que nous avons fait}
% Consigne : en langage courant, sans nom de méthode ni de package. Par exemple
% « nous avons comparé les ventes de 400 magasins » plutôt que « régression
% linéaire multiple ». Précisez la source et la taille des données.
\begin{itemize}
  \item \textbf{Données :} [À compléter : quoi, combien, quelle période, quelle
        source ; réelles ou simulées].
  \item \textbf{Démarche :} [À compléter : ce que nous avons comparé, mesuré ou
        prédit, en une ou deux phrases].
  \item \textbf{Fiabilité :} [À compléter : par exemple « vérifié sur des clients
        non utilisés pour construire le modèle »].
\end{itemize}

% ----------------------------------------------------------------------------
\section{Les principaux résultats}
% Consigne : 3 à 4 résultats, du plus important au moins important, chacun
% chiffré et traduit en conséquence pour l'entreprise.
\begin{enumerate}
  \item \textbf{[Résultat 1].} [À compléter : chiffre clé et ce qu'il signifie
        pour l'entreprise.]
  \item \textbf{[Résultat 2].} [À compléter.]
  \item \textbf{[Résultat 3].} [À compléter.]
\end{enumerate}

% Consigne : UN graphique, celui qui porte le message principal. Son titre
% énonce le message ; les axes ont des unités ; une phrase le commente.
\begin{figure}[h]
  \centering
  \IfFileExists{figures/figure-cle.pdf}{%
    \includegraphics[width=0.85\linewidth]{figures/figure-cle.pdf}%
  }{%
    \fbox{\parbox[c][4cm][c]{0.8\linewidth}{\centering\color{gray}
      Graphique principal à exporter depuis R :\\[0.3em]
      \texttt{ggsave("figures/figure-cle.pdf", width = 7, height = 3.5)}}}%
  }
  \caption{[À compléter : le message du graphique en une phrase.]}
\end{figure}

% ----------------------------------------------------------------------------
\section{Nos recommandations}
% Consigne : des actions concrètes et réalisables, avec leur effet attendu
% (chiffré si possible), leur coût ou leur difficulté, et une priorité.
\begin{center}
\rowcolors{2}{fondencadre}{white}
\begin{tabularx}{\linewidth}{>{\raggedright\arraybackslash}X
                             >{\raggedright\arraybackslash}p{4.2cm}
                             >{\centering\arraybackslash}p{1.8cm}}
  \toprule
  \textbf{Action recommandée} & \textbf{Effet attendu} & \textbf{Priorité} \\
  \midrule
  % Ne commencez pas une ligne du tableau par « [ » : LaTeX le lirait comme
  % un argument de la commande précédente (\\ ou \midrule).
  Action 1 : à compléter & Effet attendu : à compléter & Haute \\
  Action 2 : à compléter & Effet attendu : à compléter & Moyenne \\
  Action 3 : à compléter & Effet attendu : à compléter & Basse \\
  \bottomrule
\end{tabularx}
\end{center}

% ----------------------------------------------------------------------------
\section{Limites et prochaines étapes}
% Consigne : soyez honnêtes sur ce que l'analyse ne permet pas de conclure
% (données, association et non causalité...) et proposez la suite : données à
% collecter, test A/B pour confirmer un effet, suivi d'indicateurs.
\begin{itemize}
  \item \textbf{Limites :} [À compléter.]
  \item \textbf{Prochaines étapes :} [À compléter.]
\end{itemize}

\vfill
{\small\color{gray} Le détail des analyses figure dans le rapport technique du
groupe.}

\end{document}
